\documentclass[12pt,a4paper]{article}

% Encoding & Language
\usepackage[utf8]{inputenc}
\usepackage[T1]{fontenc}
\usepackage{babel}

% Layout
\usepackage[a4paper, top=2.5cm, bottom=2.5cm, left=3cm, right=2.5cm, headheight=15pt]{geometry}
\usepackage{setspace}
\onehalfspacing

% Fonts


% Graphics & Colors
\usepackage{graphicx}
\usepackage[table,xcdraw]{xcolor}
\definecolor{iselred}{RGB}{180,0,0}
\definecolor{codebg}{RGB}{245,245,245}
\definecolor{lightgray}{RGB}{230,230,230}

% Links
\usepackage[hidelinks, colorlinks=true, linkcolor=iselred, urlcolor=iselred, citecolor=iselred]{hyperref}

% Directory tree


% Code listings
\usepackage{listings}
\lstset{
  backgroundcolor=\color{codebg},
  basicstyle=\ttfamily\small,
  breaklines=true,
  frame=single,
  framesep=4pt,
  rulecolor=\color{lightgray},
  xleftmargin=10pt,
  xrightmargin=10pt,
  showstringspaces=false,
  tabsize=2,
}

% Tables
\usepackage{booktabs}
\usepackage{tabularx}
\usepackage{array}

% Headers & Footers
\usepackage{fancyhdr}
\pagestyle{fancy}
\fancyhf{}
\fancyhead[L]{\small Diário LX}
\fancyhead[R]{\small Organização do Projeto}
\fancyfoot[C]{\thepage}
\renewcommand{\headrulewidth}{0.4pt}

% Section styling
\usepackage{titlesec}
\titleformat{\section}{\large\bfseries\color{iselred}}{}{0em}{\thesection\quad}[\titlerule]
\titleformat{\subsection}{\normalsize\bfseries}{}{0em}{\thesubsection\quad}
\titleformat{\subsubsection}{\normalsize\itshape}{}{0em}{\thesubsubsection\quad}

% Misc
\usepackage{enumitem}
\usepackage{parskip}
\setlength{\parskip}{6pt}

% --------------------------------------------------
\begin{document}

% ===== CAPA =====
\begin{titlepage}
  \centering
  \vspace*{1cm}

  {\Large Instituto Superior de Engenharia de Lisboa}\\[0.3cm]
  {\large Licenciatura em Engenharia Informática e de Computadores}\\[0.2cm]
  {\large Projeto e Seminário 2025/2026}

  \vspace{2cm}
  {\LARGE \textbf{Organização do Projeto}}\\[0.5cm]
  {\Huge \textbf{\textcolor{iselred}{Diário LX}}}\\[0.3cm]
  {\large Plataforma de Jornalismo Digital}

  \vspace{2.5cm}
  \begin{tabular}{rl}
    \textbf{Estudantes:}   & Jessé Alencar, n.º 51745, \href{mailto:a51745@alunos.isel.pt}{a51745@alunos.isel.pt}   \\[0.2cm]
                           & Julianne Lobato, n.º 51191, \href{mailto:a51191@alunos.isel.pt}{a51191@alunos.isel.pt} \\[0.5cm]
    \textbf{Orientadores:} & Artur Ferreira, \href{mailto:artur.ferreira@isel.pt}{artur.ferreira@isel.pt} — ISEL    \\[0.2cm]
                           & Filipe Freitas, \href{mailto:filipe.freitas@isel.pt}{filipe.freitas@isel.pt} — ISEL    \\[0.2cm]
                           & Fátima Cardoso, \href{mailto:mlcardoso@escs.ipl.pt}{mlcardoso@escs.ipl.pt} — ESCS      \\
  \end{tabular}

  \vspace{2.5cm}
  {\large Junho de 2026}

  \vfill
\end{titlepage}

% ===== 1. REPOSITÓRIO =====
\section{Repositório do Projeto}

Toda a implementação e documentação do projeto
\textbf{Diário LX} estão disponíveis publicamente no \textit{GitHub}:

\begin{center}
  \href{https://github.com/jfmalencar/DiarioLX}{\texttt{https://github.com/jfmalencar/DiarioLX}}
\end{center}

% ===== 2. ORGANIZAÇÃO =====
\section{Organização do Projeto}

O projeto está organizado num único repositório. O código fonte encontra-se na pasta \texttt{code/}, dividida em dois diretórios principais: o \textit{backend} (JVM/Kotlin) e o \textit{frontend} (React/TypeScript).
\subsection{\textit{Backend} — \texttt{code/jvm/}}

\begin{lstlisting}
code/jvm/
|-- modules/
|   |-- domain/           % Entidades e regras de negocio
|   |-- http/             % Controllers, DTOs e autenticacao
|   |-- services/         % Logica de negocio e validacao
|   |-- repository/       % Interfaces de acesso a dados
|   |-- repository-jdbi/  % Implementacoes JDBI + SQL
|   |-- storage/          % Interface de armazenamento
|   |-- storage-s3/       % Implementacao S3 / SeaweedFS
|   `-- host/             % Ponto de entrada e configuracao
|-- build.gradle.kts
|-- settings.gradle.kts
|-- seaweed-config.json
`-- docker-compose.yml
\end{lstlisting}

\subsection{\textit{Frontend} — \texttt{code/js/}}

\begin{lstlisting}
code/js/
|-- src/
|   |-- app/
|   |   `-- providers/    % Bootstrap, autenticacao e i18n
|   |-- assets/           % Logotipo, icones e traducoes
|   |-- config/           % Variaveis de ambiente
|   |-- layouts/          % Layouts partilhados
|   |-- modules/
|   |   |-- backoffice/   % Interface de gestao editorial
|   |   `-- public/       % Interface publica do jornal
|   |-- shared/
|   |   |-- components/   % Componentes reutilizaveis
|   |   |-- hooks/        % Hooks de acesso ao servicos
|   |   |-- services/     % Interfaces e implementacoes
|   |   |-- types/        % Tipos e interfaces partilhados
|   |   `-- utils/        % Funcoes utilitarias
|   `-- tests/            % Testes Playwright
|-- index.html
|-- index.tsx
`-- package.json
\end{lstlisting}

\section{Implantação e Utilização}

\subsection{Pré-requisitos}

\begin{itemize}[leftmargin=1.5cm]
  \item \textbf{Docker} e \textbf{Docker Compose} instalados.
  \item \textbf{Java 17+} instalado (necessário para executar o \textit{Gradle Wrapper}).
  \item Acesso à linha de comandos (terminal).
\end{itemize}

\subsection{Modos de Execução}

O projecto suporta dois modos de execução distintos:

\begin{description}
  \item[Desenvolvimento local] O \textit{backend} é executado directamente pelo
        \textit{IntelliJ IDEA} e o \textit{frontend} pelo servidor de desenvolvimento \textit{Vite}
        (\texttt{npm run dev}). As variáveis de ambiente são carregadas
        a partir do ficheiro \texttt{.env} na raiz de \texttt{code/jvm/}.
        Apenas a base de dados e o \textit{SeaweedFS} são executados via \textit{Docker
          Compose}:

        \begin{lstlisting}[language=bash]
cd code/jvm
docker compose up -d
\end{lstlisting}

  \item[\textit{Staging}] Todos os serviços — \textit{backend}, base de dados,
        \textit{SeaweedFS} e a imagem \textit{web} (\textit{frontend} servido por
        \textit{Nginx}, que actua também como \textit{reverse proxy}) — são executados em
        \textit{containers} \textit{Docker}. As variáveis de ambiente são carregadas a partir do
        ficheiro \texttt{.env.staging}. Este é o modo descrito nas
        secções seguintes, usando as \textit{tasks} Gradle \texttt{buildImageAll}
        e \texttt{allUp}.
\end{description}
\subsection{Configuração Inicial}

Antes de executar qualquer comando, é necessário colocar dois ficheiros de configuração na pasta \texttt{code/jvm/}:

\subsubsection{Ficheiro \texttt{.env.staging}}

Criar o ficheiro \texttt{code/jvm/.env.staging} com o seguinte conteúdo:

\begin{lstlisting}
PORT=8180

S3_PUBLIC_ENDPOINT=http://localhost:${PORT}
S3_ENDPOINT=http://diariolx-seaweed-s3:8333
S3_REGION=us-east-1
S3_ACCESS_KEY=DEV-ACCESS-KEY
S3_SECRET_KEY=DEV-SECRET-KEY
S3_BUCKET=media
S3_PATH_STYLE_ACCESS=true

SCALAR_PATH=/docs

JWT_SECRET=QBqCNu2Dusr161yq7foRF2U6MLjKSG4oG8j7T7RNbk4=
JWT_ACCESS_EXPIRATION_MS=600000
JWT_REFRESH_EXPIRATION_MS=604800000

POSTGRES_USER=dbuser
POSTGRES_PASSWORD=changeit
POSTGRES_DB=db
POSTGRES_HOST=diariolx-postgres-test
POSTGRES_PORT=5432

DB_URL=jdbc:postgresql://${POSTGRES_HOST}:${POSTGRES_PORT}/${POSTGRES_DB}\
?user=${POSTGRES_USER}&password=${POSTGRES_PASSWORD}
\end{lstlisting}

\subsubsection{Ficheiro \texttt{seaweed-config.json}}

Alterar o ficheiro \texttt{code/jvm/seaweed-config.json} com o seguinte conteúdo:

\begin{lstlisting}
{
  "identities": [
    {
      "name": "anonymous",
      "actions": ["Read"],
      "resources": ["buckets/media/*"]
    },
    {
      "name": "dev-user",
      "credentials": [
        {
          "accessKey": "DEV-ACCESS-KEY",
          "secretKey": "DEV-SECRET-KEY"
        }
      ],
      "actions": ["Read", "Write", "List", "Tagging", "Admin"],
      "resources": ["buckets/media/*"]
    }
  ]
}
\end{lstlisting}

Este ficheiro configura as permissões do \textit{SeaweedFS} — o sistema de armazenamento de ficheiros multimédia. O utilizador \texttt{anonymous} tem permissão de leitura pública, enquanto o \texttt{dev-user} tem acesso total para operações do \textit{backend}.

\subsection{Construção e Arranque}

Com os ficheiros de configuração em lugar, navegar para a pasta \texttt{code/jvm/} e executar os seguintes comandos pela ordem indicada.

\subsubsection{Passo 1 — Construir as imagens \textit{Docker}}

\begin{lstlisting}[language=bash]
cd code/jvm
./gradlew buildImageAll
\end{lstlisting}

Este comando compila o projecto e constrói todas as imagens \textit{Docker} necessárias (\textit{backend}, \textit{frontend} e infraestrutura).

\subsubsection{Passo 2 — Iniciar todos os serviços}

\begin{lstlisting}[language=bash]
./gradlew allUp
\end{lstlisting}

Este comando inicia todos os serviços orquestrados via \textit{Docker Compose}. Após a conclusão, a aplicação estará disponível na porta \texttt{8180}.

\subsection{Resumo dos \textit{Endpoints}}

Após o arranque, todos os serviços são acessíveis através de uma única porta:

\begin{center}
  \begin{tabularx}{0.85\textwidth}{@{}lXl@{}}
    \toprule
    \textbf{\textit{Endpoint}} & \textbf{Descrição}                                 & \textbf{URL}                      \\
    \midrule
    \textit{Frontend}          & Interface web (site público e \textit{backoffice}) & \texttt{localhost:8180}           \\
    API REST                   & \textit{Backend} Spring Boot                       & \texttt{localhost:8180/api}       \\
    Documentação               & Scalar (especificação \textit{OpenAPI})            & \texttt{localhost:8180/docs}      \\
    \textit{Media}             & Ficheiros multimédia (\textit{SeaweedFS})          & \texttt{localhost:8180/media/...} \\
    \bottomrule
  \end{tabularx}
\end{center}

\noindent\colorbox{codebg}{\parbox{\dimexpr\linewidth-2\fboxsep}{%
    \small\textbf{Nota:} O \textit{Nginx} actua como \textit{reverse proxy}, encaminhando todos os
    pedidos através da porta \texttt{8180}. Não é necessário expor portas individuais
    de cada serviço.
  }}

\subsection{Convite Inicial de Administrador}

Na primeira execução, é inserido automaticamente na base de dados um convite de administrador com o código:

\begin{lstlisting}
SUPER-ADMIN-INVITE
\end{lstlisting}

Aceder a \texttt{localhost:8180/backoffice/register} e
usar este código para criar uma conta com o papel de \texttt{ADMIN}.

% ===== 4. TECNOLOGIAS =====
\section{Tecnologias Utilizadas}

\begin{center}
  \begin{tabularx}{0.85\textwidth}{@{}lX@{}}
    \toprule
    \textbf{Tecnologia}  & \textbf{Utilização}                                         \\
    \midrule
    Kotlin + Spring Boot & \textit{Backend} e API REST                                 \\
    JDBI                 & Acesso à base de dados                                      \\
    PostgreSQL           & Persistência de dados estruturados                          \\
    SeaweedFS (S3)       & Armazenamento de ficheiros multimédia                       \\
    React + TypeScript   & Interface \textit{web} (\textit{backoffice} e site público) \\
    Vite                 & \textit{Bundler} e servidor de desenvolvimento              \\
    Playwright           & Testes \textit{end-to-end} do \textit{frontend}             \\
    Docker + Compose     & \textit{Containerização} e orquestração de serviços         \\
    Nginx                & \textit{Reverse proxy} (implantação)                        \\
    \bottomrule
  \end{tabularx}
\end{center}

\end{document}